\documentclass[../../main.tex]{subfiles}
\begin{document}

\begin{center}
  {\fontsize{48pt}{58pt}\selectfont \textbf{1969}\normalsize\textbf{年度}}
\end{center}

\begin{table}[htb]
  \centering
  \caption{出題テーマと難易度}
  \begin{tabular}{|c|p{9cm}|c|}
    \hline
    \textbf{問題番号} & \textbf{テーマ} & \textbf{難易度} \\
    \hline
    第1問 & 不等式の証明 & \\
    \hline
    第2問 & 空間図形の計量 & \\
    \hline
    第3問 & 軌跡と領域の図示 & \\
    \hline
    第4問 & 数列と図形の面積 & \\
    \hline
    第5問 & 方程式の実根の存在証明 & \\
    \hline
    第6問 & 微分法と大小比較 & \\
    \hline
  \end{tabular}
\end{table}

\begin{table}[htb]
  \centering
  \caption{解答}
  \renewcommand{\arraystretch}{1.4}
  \begin{tabular}{|c|c|p{9cm}|}
    \hline
    \textbf{問題番号} & \textbf{小問} & \textbf{答え} \\
    \hline
    第1問 & - & 証明問題（$x^2-y^2-z^2<0$を示す） \\
    \hline
    第2問 & (1) & $DE=\sqrt{\dfrac{433}{13}}$ \\
    \cline{2-3}
    & (2) & $AF=30\sqrt{\dfrac{3}{433}}$ \\
    \hline
    第3問 & (1) & $v \ge \dfrac{1}{2}(-u^2+2u+1)$，$v=-u^2+2u+a^2$上 \\
    \cline{2-3}
    & (2) & 図示（斜線部の領域，境界含む） \\
    \hline
    第4問 & - & $\displaystyle\lim_{n\to\infty}a_n=1$ \\
    \hline
    第5問 & - & 証明問題 \\
    \hline
    第6問 & - & $x^2e^{y/x}>y^2e^{x/y}$ \\
    \hline
  \end{tabular}
\end{table}

\paragraph{試験形態}
6問．試験時間は要確認．

\end{document}
